\chapter{Анатомія токамака: JET}\label{ch:jet}

% Рамка «Модель проєкту» (числа нашої реконструкції). Визначено тут через \ProvideTColorBox,
% щоб не редагувати preamble.tex; координатор може перенести її в преамбулу.
\ProvideTColorBox{projmodel}{O{}}{colback=violet!4,colframe=violet!55!black,
  title={Модель проєкту\ifx&#1&\else{} --- #1\fi}}

У попередніх розділах токамак був переважно рівнянням: $\psi(R,Z)$, $q(\psi)$, метрика. Тут розберемо реальну машину шар за шаром, від плазми до фундаменту, і з'ясуємо, чому кожен шар має саме таку форму. Приклад --- JET (Joint European Torus) у двох станах. Перший --- проєкт 1975~р. (EUR~5516e), де більшість розмірів надруковано на кресленнях~\cite{JET1975}. Другий --- машина 1983~р., як її побудовано до першої плазми 25.06.1983~\cite{JETJU1983}. За цими даними проєкт має власну 3D-реконструкцію~\cite{ProjJET} (рис.~\ref{fig:04j-section}, \ref{fig:04j-cutaway}). У джерелах немає жодного рисунка, який би ми відтворювали: усі зображення тут --- наші рендери й графіки, а з документів узято лише числа.

\section{Чому саме така форма}\label{sec:04j-shape}

\paragraph{Мале аспектне відношення.} Головну величину JET обрано першою: струм плазми $\Ip\sim\SI{3}{МА}$. За тодішніми даними добуток $n\tauE$ зростав з $\Ip$, а за такого струму орбіти значної частини $\alpha$-частинок лишаються в плазмі, поки вони віддають енергію~\cite[с.~67]{JET1975}. Далі постає питання ціни. У наближенні кругового перерізу й великого аспектного відношення автори проєкту записують три співвідношення~\cite[с.~74]{JET1975}:
\begin{equation}\label{eq:04j-Ic}
  \ITF=q\Big(\frac{R_0}{a}\Big)^{2}\Ip,\qquad
  \Ip=\frac{2\pi a\Bpol}{\mu_0},\qquad
  \ITF=\frac{2\pi}{\mu_0}\Bleg\,(R_0-a-e_1),
\end{equation}
де $\ITF$ --- повний струм котушок TF, $\Bleg$ --- поле на внутрішній нозі котушки (у джерелі $B_{\mathrm{MAX}}$), обмежене механічними напруженнями, $e_1$ --- товщина камери. Перша рівність --- це просто $q(a)$ з~\eqref{eq:01-qcyl}, записане через $\Btor=\mu_0\ITF/(2\pi R_0)$. Звідси видно: за заданих $\Ip$ і $q$ струм TF, а отже й вартість магніту, зростає як $(R_0/a)^2$. Тому $R_0/a$ беруть якомога меншим. Знизу його обмежує трансформатор: потік через центральне осердя має встигнути навести струм. Поєднання~\eqref{eq:04j-Ic} з рівнянням трансформатора дає мінімум $R_0/a\approx2{,}2$; для JET обрано 2,37~\cite[с.~75]{JET1975}. Перевірка: $q=2{,}8$, $\Ip=\SI{2,6}{МА}$, $R_0/a=2{,}37$ дають $\ITF=\SI{40,9}{МА}$ проти проєктних \SI{41}{МА}~\cite[с.~83]{JET1975}. Поле на осі менше за поле на нозі: $\Btor=\Bleg(1-a/R_0-e_1/R_0)$, тож при $\Bleg=\SI{6}{Тл}$ маємо $\Btor\approx\SI{3}{Тл}$~\cite[с.~75]{JET1975}. Розрахунок повної вартості дав той самий оптимум, $R_0/a\approx2{,}4$ при \SI{3}{Тл}~\cite[с.~68]{JET1975}.

\paragraph{D-подібна котушка.} Аспектне відношення можна зменшувати лише з тонкою котушкою TF, без масивного корпусу~\cite[с.~68]{JET1975}. Така котушка повинна сама тримати магнітний тиск. Кругова котушка з цим не впорається: вона згинається, і в ізоляції між витками виникає зсув до \SI{2}{кгс/мм^2} ($\approx\SI{20}{МПа}$)~\cite[с.~333]{JET1975}. Епоксидне зчеплення з міддю такого не витримує. Вихід --- надати котушці форму, за якої згину немає взагалі.

\begin{outofcorpus}[котушка чистого натягу («D принстонського типу»)]
Нехай $N$ тонких котушок зі струмом $I_1$ кожна ($\ITF=NI_1$) утворюють поле $\Btor=\mu_0\ITF/(2\pi R)$. Сама обмотка перебуває в середньому полі $\Btor/2$ (поле спадає до нуля поперек обмотки), тому на одиницю довжини діє нормальна сила назовні
\begin{equation}\label{eq:04j-f}
  f(R)=\frac{\mu_0\ITF{}I_1}{4\pi R}.
\end{equation}
Гнучка нитка під нормальним навантаженням має сталий натяг $\Tcoil$ (дотичної сили немає), а рівновага елемента дає $\Tcoil/\rho=f$, де $\rho$ --- радіус кривини. Звідси
\begin{equation}\label{eq:04j-rho}
  \rho(R)=\kD{}R,\qquad \Tcoil=\frac{\mu_0\ITF{}I_1}{4\pi}\,\kD,\qquad \kD=\tfrac12\ln\frac{R_2}{R_1}.
\end{equation}
Сталу $\kD$ знаходимо з балансу вертикальних сил на верхній половині контуру між $R_1$ (верх прямої внутрішньої ноги) і $R_2$ (зовнішній екватор). В обох точках дотична вертикальна, тож натяг на двох розрізах тягне вниз з силою $2\Tcoil$, а вертикальна складова нормальної сили дорівнює $\int f\,n_Z\,\dd\ell=\int_{R_1}^{R_2}f\,\dd R$. Отже, $2\Tcoil=\mu_0\ITF{}I_1\ln(R_2/R_1)/(4\pi)$, тобто $\Tcoil=\mu_0NI_1^2\ln(R_2/R_1)/(8\pi)$ на одну котушку. Радіус кривини зростає з $R$: біля внутрішньої ноги контур вигинається круто, а зовнішня частина --- полога дуга. Разом із прямою внутрішньою ногою це й дає літеру D. На прямій нозі ($\rho=\infty$) нормальну силу натяг не врівноважує; її сумарно передають на центральну опору. Це доцентрова сила, розглянута нижче.
\end{outofcorpus}

Саме цим критерієм --- сила на елемент $\Delta\ell$ стала, якщо $\Delta\ell\propto\rho$, тобто $\rho\propto R$ --- автори перевіряли форму котушки JET. Відправною точкою був розв'язок для суцільної обмотки. Потім форму уточнювали ітераціями, щоб урахувати широкі проміжки між 32 котушками на зовнішньому периметрі; результат виявився на $\approx\SI{5}{см}$ пласкішим~\cite[с.~333--334]{JET1975}. Розрахунок методом скінченних елементів показав, що натяг уздовж контуру змінюється лише на $\pm12\,\%$, а зсув не перевищує \SI{0,2}{кгс/мм^2}. Це вдесятеро менше, ніж у кругової котушки~\cite[с.~336]{JET1975}.

\begin{projmodel}[натяг котушки JET]
Осьова лінія обмотки нашої моделі проходить від $R_1=\SI{1,304}{м}$ до $R_2=\SI{4,790}{м}$~\cite{ProjJET}. За~\eqref{eq:04j-rho} при $\ITF=\SI{41}{МА}$, $I_1=\SI{1,28}{МА}$ отримуємо $\Tcoil=\SI{3,42}{МН}$. На внутрішній нозі мідь має переріз $24\times\SI{2700}{мм^2}$~\cite[с.~327]{JET1975}, тож напруження розтягу $\approx\SI{53}{МПа}$ (5,4~кгс/мм$^2$). Проєктний розрахунок дає 5,9~кгс/мм$^2$ на зменшеному перерізі~\cite[с.~339]{JET1975}. Масштабування до \SI{51}{МА} дає $\approx\SI{5,3}{МН}$ ($\approx540$~тс), що узгоджується з оцінкою «до 600~т на котушку»~\cite[с.~40]{Wesson1999}.
\end{projmodel}

\paragraph{D-подібна плазма.} Кругла плазма в D-подібному отворі лишає порожніми «кути». За тієї самої площі поверхні тор найбільшого об'єму має переріз, близький до D, і оптимальна котушка теж близька до D~\cite[с.~70]{JET1975}. Тому плазму витягнуто до $b/a=1{,}68$ ($a=\SI{1,25}{м}$, $b=\SI{2,10}{м}$, $R_0=\SI{2,96}{м}$)~\cite[с.~83]{JET1975}, і вона майже заповнює отвір. За теоретичними оцінками проєкту витягнутість $\approx1{,}7$ дає в $1{,}7$ раза більший струм за того самого $q(0)$~\cite[с.~70]{JET1975} (розд.~\ref{ch:geometry}). Для D-перерізу проєкт оцінює
\begin{equation}\label{eq:04j-qD}
  \frac{\Ip\,q(a)}{\ITF}\approx k\,\frac{1+(b/a)^2}{2}\,\frac{(a/R_0)^2}{\bigl[1-(a/R_0)^2\bigr]^2},\qquad k=1{,}1,
\end{equation}
де $k$ взято з розрахунку пікованого струму з $\betap=0{,}9$~\cite[с.~83]{JET1975}. Звідси \SI{3,8}{МА} при $q(a)=6$. Сподівалися, що $q(0)\approx1$ відповідатиме $q(a)\approx6$~\cite[с.~71]{JET1975}. Ціна малого $R_0/a$ --- більше захоплених частинок: $\approx77\,\%$ на зовнішній поверхні JET проти 63\,\% при $R_0/a=4$~\cite[с.~71]{JET1975}. Це збігається з оцінкою $\sqrt{2\epsilon/(1+\epsilon)}$ при $\epsilon=a/R_0$.

\section{Шари від плазми назовні}\label{sec:04j-layers}

\begin{figure}[!tb]\centering
  \includegraphics[width=0.88\textwidth]{jet_section.jpg}
  \caption{Вертикальний переріз JET 1975 через плече ярма. Показано розміри з посиланнями на сторінки джерела; кольори кодують матеріали. Від осі назовні: центральна колона ярма, стек PF-1, внутрішня нога TF, камера, плазма ($\psiN$), зовнішня нога TF, PF-2--4, вертикальні плечі. Власна 3D-реконструкція~\cite{ProjJET}; побудовано за даними~\cite[с.~4, 83, 298, 378, 401, 405--406]{JET1975}.}
  \label{fig:04j-section}
\end{figure}

\paragraph{Плазма і перша стінка.} Межа плазми лежить на $R=1{,}71\ldots\SI{4,21}{м}$. Від внутрішньої стінки жорсткого сектора її відділяють \SI{5}{см}, від зовнішньої --- \SI{16}{см}~\cite[с.~83]{JET1975}. Плазма має торкатися не стінки, а лімітерів. У проєкті 1975~р. їх два види. Основні --- три кільцеві «рейки» з окремих пластин уздовж тора; пластини охолоджує потік газу, і їхнє положення можна підправляти ззовні, не відкриваючи камеру~\cite[с.~301]{JET1975}. Зовнішню рейку складали 16 молібденових пластин $100\times60\times\SI{1}{см}$, разом близько тонни Mo~\cite[с.~314]{JET1975}. Допоміжні --- полоїдальні екрани біля кожного сильфона на $R=1{,}70$ і \SI{4,33}{м}: тонкостінні сильфони найвразливіші, зокрема до утікачів~\cite[с.~83, 301]{JET1975}. Уже в проєкті матеріал лімітера вважали «експериментальним параметром», і важкий метал могли замінити на легкий~\cite[с.~301, 315]{JET1975}. У 1983~р. так і сталося (п.~\ref{sec:04j-1983}).

\paragraph{Вакуумна камера.} Форму камери диктує магніт: до котушок TF лишено лише зазори на монтаж, теплоізоляцію і рух стінок під час прогріву й навантаження. Конструкція цілком металева й зварна, як того вимагає надвисокий вакуум~\cite[с.~295]{JET1975}. Стінка --- «коробка» з двох обшивок інконелю: 20~мм на прямій внутрішній ділянці, 12~мм далі по периметру; разом з проміжком це \SI{12}{см} на екваторі й до \SI{18}{см} угорі~\cite[с.~298, 310]{JET1975}. По тору чергуються два типи елементів. Жорсткі клиноподібні сектори (їх 32, по одному на котушку TF) несуть усі сили й мають отвори портів. Між ними стоять сильфонні вставки: пара вкладених один в один сильфонів з листа \SI{1,7}{мм}, проміжок між якими має окреме відкачування на випадок течі внутрішнього сильфона~\cite[с.~295, 306]{JET1975}. Одиниця збирання --- октант із чотирьох секторів і чотирьох сильфонів; шов між октантами проходить посередині крайнього сектора, тому пошкоджений октант можна вирізати й замінити дистанційно цілим модулем~\cite[с.~297, 299]{JET1975}. Маса камери \SI{68}{т}, об'єм $\approx\SI{190}{м^3}$. Площа внутрішньої поверхні --- $\approx\SI{220}{м^2}$ для гладкого тора, $\approx450$ з гофрами сильфонів і $\approx\SI{780}{м^2}$ з екранами й портами~\cite[с.~310]{JET1975}. Цільовий базовий тиск --- $2\cdot10^{-10}$~Торр ($\approx\SI{2,7e-8}{Па}$)~\cite[с.~315]{JET1975}.

\emph{Навіщо сильфони? Тороїдальний опір.} Камера --- замкнений провідний виток навколо того самого осердя трансформатора, що й плазма. Тому напруга обходу $\Vloop$ (розд.~\ref{ch:trap}) наводить у ній паралельний струм $I_{\mathrm v}=\Vloop/\Rv$. Цей струм марно витрачає вольт-секунди, спотворює полоїдальне поле біля плазми і створює сили на камеру. Сильфони мають мале січення й довгий гофрований шлях струму, тож саме вони задають опір камери «довгим шляхом навколо тора». Його свідомо роблять якомога більшим, щоб струм ішов через газ, а не через стінки~\cite[с.~295]{JET1975}, \cite[с.~12]{JETBrochure1982}: $\Rv=\SI{0,52}{мОм}$, індуктивність $\Lv\approx\SI{2,5}{мкГн}$~\cite[с.~304, 310]{JET1975}. Це приблизно вп'ятеро більше за опір водневої плазми з $T_e=\SI{5}{еВ}$, яка заповнює всю камеру~\cite[с.~304]{JET1975}. Отже, на пробої струм іде переважно в плазму. Стала часу $\Lv/\Rv\approx\SI{4,8}{мс}$ задає, як швидко камера «перехоплює» струм під час раптового зриву; для найгіршого випадку проєкт приймав, що в камеру переходить увесь $\Ip$~\cite[с.~304]{JET1975}. У полоїдальному напрямку всі жорсткі сектори ввімкнено паралельно, тож опір там малий (\SI{0,03}{мОм}). Під час наростання $\Btor$ це дає вихрові струми й тиск на стінки сектора~\cite[с.~303--304, 310]{JET1975}.

\begin{outofcorpus}[дві оцінки для камери]
\emph{Без сильфонів.} Уявімо суцільну подвійну обшивку з інконелю ($\rhoe\approx\SI{1,0e-6}{Ом.м}$) із середньою товщиною $2\times\SI{15}{мм}$ на полоїдальному периметрі $\approx\SI{11}{м}$. Для кола довжиною $2\pi R_0$ це дає $\mathcal{R}\approx\rhoe\,2\pi R_0/(0{,}33~\text{м}^2)\approx\SI{0,06}{мОм}$, тобто вдев'ятеро менше, ніж з сильфонами. \emph{Плазма на пробої.} Спітцерівський опір $\eta\approx5{,}2\cdot10^{-5}\lnL/T_e^{3/2}$~Ом$\cdot$м ($T_e$ в еВ) при $T_e=\SI{5}{еВ}$, $\lnL\approx10$ і перерізі камери $\approx\SI{9}{м^2}$ дає $\approx\SI{0,1}{мОм}$. Відношення $\approx5$ збігається з наведеним у проєкті.
\end{outofcorpus}

\emph{Прогрів.} Щоб досягти надвисокого вакууму, камеру прогрівають до \SI{500}{\degreeCelsius} гарячим інертним газом, який проганяють через простір між обшивками. Ззовні камеру вкрито теплоізоляцією товщиною \SI{2,5}{см}~\cite[с.~299, 310, 320]{JET1975}. Теплове розширення становить $\approx\SI{30}{мм}$~\cite[с.~312]{JET1975}, і всі зазори до котушок розраховано з цим запасом (перевірте це в задачі~2).

\paragraph{Котушки TF.} Магніт утворюють 32 D-подібні мідні котушки з водяним охолодженням. Кожна котушка складається з двох «млинців» по 12 витків~\cite[с.~323--324]{JET1975}; габарит $5{,}68\times\SI{3,86}{м}$, маса \SI{12}{т}~\cite[с.~327]{JET1975}, увесь магніт важить $\approx\SI{400}{т}$~\cite[с.~323]{JET1975}. Внутрішні ноги мають клиноподібний переріз і змикаються навколо центрального стека PF. Тому мідь усередині має лише 70\,\% перерізу зовнішньої частини (2700 проти \SI{3900}{мм^2} на виток)~\cite[с.~324, 327]{JET1975}. Струм \SI{53,4}{кА} на виток ($\ITF=\SI{41}{МА}$) дає $\Bzero=\SI{2,77}{Тл}$ на $R=\SI{2,96}{м}$. Розширений режим --- \SI{66,4}{кА} (\SI{51}{МА}, \SI{3,45}{Тл}). Запасена енергія становить 0,94 і \SI{1,45}{ГДж}, повна індуктивність \SI{0,66}{Гн}~\cite[с.~328]{JET1975}. Кожну котушку тягне до осі доцентрова сила $1216$~тс ($\approx\SI{11,9}{МН}$), а в розширеному режимі --- $1880$~тс. Цю силу сприймає клиноподібний «носик» внутрішньої ноги, що впирається в сусідні котушки та центральну опору~\cite[с.~335, 339]{JET1975}.

\paragraph{Механічна структура.} Власне поле TF навантажує котушку лише в її площині. Інша справа --- полоїдальне поле: воно перетинає котушку, і сила $I_1\,\dd\vect{l}\times\Bvec_{\mathrm p}$ напрямлена вже по $\tor$. Оскільки $\BR$ вище й нижче екватора має різні знаки, ця сила тягне верхню й нижню половини котушки в протилежні боки по $\tor$, а від усіх 32 котушок складається крутний момент навколо вертикальної осі~\cite[с.~359]{JET1975}. Його сприймають тонка оболонка з клинових блоків між котушками, верхнє й нижнє кільця ($\varnothing\,\SI{7,32}{м}$, висота \SI{0,9}{м}, по 8 секторів) і внутрішній циліндр з 8 секторів нержавіючої сталі завтовшки в середньому \SI{4}{см}~\cite[с.~364, 374--378]{JET1975}. У збудованій машині оболонку й кільця відлито з немагнітного аустенітного високоміцного чавуну (з кулястим графітом), а внутрішній циліндр зроблено зі сталі 304 з 32 пазами під ноги котушок. Структура важить $\approx\SI{470}{т}$ і розрахована на момент \num{20000}~тс$\cdot$м ($\approx\SI{2e8}{Н.м}$) при закручуванні оболонки не більш ніж на \SI{4}{мм}~\cite[с.~16--17]{JETBrochure1982}.

\paragraph{Котушки PF.} Полоїдальну систему поставлено \emph{ззовні} котушок TF. Так дві системи не зачеплені топологічно, і будь-яку з них можна розібрати, не чіпаючи іншої~\cite[с.~72]{JET1975}. 16 мідних котушок чотирьох типів (табл.~\ref{tab:04j-pf}) утворюють первинну обмотку трансформатора (PF-1) і котушки рівноваги та форми (PF-2--4)~\cite[с.~385, 399]{JET1975}. Для швидких процесів обмотки діють як паралельні контури: різкий зсув плазми змінює їхнє потокозчеплення, і наведені струми (правило Ленца) тягнуть її назад, як суцільна мідна оболонка~\cite[с.~385]{JET1975} (розд.~\ref{ch:stability}).

\begin{table}[!tb]\centering\small\setlength{\tabcolsep}{5pt}
\caption{Котушки PF JET у проєкті 1975~р.~\cite[с.~401]{JET1975}}\label{tab:04j-pf}
\begin{tabular}{@{}lcccc@{}}
\toprule
 & PF-1 (внутр.) & PF-2 (внутр.) & PF-3 (зовн.) & PF-4 (зовн.) \\
\midrule
Кількість & 10 & 2 & 2 & 2 \\
Середній діаметр, м & 1,816 & 4,3 & 7,94 & 10,486 \\
Відстань від екватора $\pm Z$, м & --- (стек) & 3,65 & 2,86 & 1,075 \\
Витків (послідовно) на котушку & 63 & 48 & 120 & 96 \\
Макс.\ струм у провіднику, кА & 34 & 34 & 11,8 & 13 \\
Маса міді на котушку, т & 5,7 & 10 & 20,2 & 35,9 \\
\bottomrule
\end{tabular}
\end{table}

\paragraph{Залізне осердя трансформатора.} Струм у плазмі наводиться трансформатором: $\Vloop=-\dd\Phi/\dd t$, де $\Phi$ --- потік, зчеплений з плазмовим витком. Ярмо --- це 8 вертикальних плечей, 8 верхніх і 8 нижніх радіальних плечей, два центральні блоки й центральна колона, разом 27 основних деталей і \SI{2263}{т}~\cite[с.~415]{JET1975}. Навіщо залізо, якщо центральна колона все одно насичується? Зовнішній контур з ненасиченого заліза ($B<\SI{1,9}{Тл}$) дає потоку короткий шлях повернення. Через це машині потрібно лише $\approx2/3$ струму намагнічування і $\approx1/2$ енергії порівняно з еквівалентною машиною без заліза~\cite[с.~399, 411]{JET1975}. Крім того, залізо притримує поле розсіювання навколо машини~\cite[с.~385]{JET1975}. Це важливо для інжекторів: уже $\sim$\SI{1}{мТл} (10~Гс) заважає їхній роботі, а біля зовнішніх плечей очікували $\sim$\SI{20}{мТл}~\cite[с.~419]{JET1975}. Ціна малого $R_0/a$ --- брак місця в центрі: на колону лишається малий переріз, і в базовому режимі залізо там насичене до \SI{6,6}{Тл}~\cite[с.~385, 412]{JET1975}. Другий прийом --- використати петлю намагнічування з обох боків: перед імпульсом осердя заздалегідь намагнічують «від'ємно», і під час наростання струму потік проходить від $-15$ до $+\SI{10}{Вб}$, тобто майже вдвічі більший перепад за тих самих струмів котушок~\cite[с.~411--412]{JET1975}. Перепад \SI{25}{Вб} і дорівнює проєктному запасу \SI{25}{В.с}~\cite[с.~83]{JET1975}.

\paragraph{Порти й доступ.} Кожен октант має великий горизонтальний порт в екваторіальній площині з отвором $0{,}46\times\SI{0,96}{м}$. Через нього пучок нейтралів діаметром до \SI{280}{мм} можна ввести під кутом $45^\circ$ до осі~\cite[с.~301]{JET1975}. Вертикальні наскрізні порти мають розміри $14\times83$ і $7\times\SI{28}{см}$~\cite[с.~77]{JET1975}. На чотирьох із восьми горизонтальних портів стоять насосні блоки~\cite[с.~315]{JET1975}. Концепція нагріву 1975~р. передбачала шість квазітангенціальних інжекторів на один порт~\cite[с.~426]{JET1975} з енергією $\ge\SI{80}{кеВ}$ (H) і нарощуванням потужності до $\sim\SI{25}{МВт}$~\cite[с.~418]{JET1975}. Уся машина має габарит $\varnothing\,14{,}8\times\SI{11,5}{м}$~\cite[с.~4]{JET1975} і масу близько \SI{3000}{т}, яку несуть нижні плечі ярма~\cite[с.~414]{JET1975}.

\begin{figure}[!tb]\centering
  \begin{subfigure}{0.495\textwidth}\centering
    \includegraphics[width=\textwidth,trim=0 40 0 130,clip]{jet_cutaway1975.jpg}\caption{}
  \end{subfigure}\hfill
  \begin{subfigure}{0.495\textwidth}\centering
    \includegraphics[width=\textwidth,trim=0 40 0 130,clip]{jet_exploded.jpg}\caption{}
  \end{subfigure}
  \caption{JET 1975: (a)~розріз на 1,5 октанта (ярмо, PF-1--4, котушки TF, оболонка, камера із сильфонами, плазма з силовими лініями); (b)~схематично рознесений октант --- одиниця збирання й дистанційної заміни (ярмо, кільце, оболонка, котушки TF і сектор камери зсунуто назовні). Людська фігура --- \SI{1,8}{м}. Власна 3D-реконструкція~\cite{ProjJET}; побудовано за даними~\cite[с.~4, 83, 295--299, 325, 378, 401, 415]{JET1975}.}
  \label{fig:04j-cutaway}
\end{figure}

\section{Гофрування поля 32 котушок}\label{sec:04j-ripple}

Скінченне число котушок робить $|\Bvec|$ періодичним по $\tor$ з періодом $2\pi/32$. Чим це загрожує швидким частинкам, розглянуто в п.~\ref{sec:02-ripple} (розд.~\ref{ch:particles}). Число котушок JET обрано саме з умови однорідності поля~\cite[с.~13]{JETBrochure1982}, а ще на початку проєктування гофрування рахували для кількох розкладок~\cite[с.~329]{JET1975}. \emph{Пастка в означеннях.} Сучасна глибина гофрування --- $\dripple=(B_{\max}-B_{\min})/(B_{\max}+B_{\min})$, де $B_{\max}$ береться в площині котушки, а $B_{\min}$ --- посередині між котушками. Але проєкт JET у підписі до графіка гофрування визначає $\eripple=2(B_{\max}-B_{\min})/(B_{\max}+B_{\min})=2\dripple$~\cite[с.~330]{JET1975}. Отже, проєктне «гофрування 3,6\,\%» на краю плазми~\cite[с.~83]{JET1975} означає $\dripple\approx1{,}8\,\%$.

\begin{projmodel}[гофрування за Біо--Саваром]
Кожну з 32 котушок представлено ниткою вздовж осьової лінії обмотки: оцифрований контур отвору~\cite[с.~325]{JET1975}, зсунутий на половину товщини ноги (\SI{0,186}{м}). Струм --- \SI{41}{МА}. Результат~\cite{ProjJET} (рис.~\ref{fig:04j-ripple}):
\begin{itemize}
\item $\eripple(\SI{4,21}{м})=3{,}67\,\%$ (з $3\times3$ нитками на котушку 3,74\,\%) проти проєктних 3,6\,\%; на осі $\dripple\approx3\cdot10^{-7}$;
\item уздовж 30 точок межі з таблиці проєкту~\cite[с.~332]{JET1975} відношення «наше/джерело» становить у середньому 0,97 (0,90--1,10);
\item незалежна перевірка: $B(R_0)=\SI{2,770}{Тл}$ при \SI{41}{МА} і \SI{3,446}{Тл} при \SI{51}{МА}, як у проєкті (2,77 і 3,45).
\end{itemize}
Головна невизначеність --- сама форма котушки. Радіуси D на кресленні не надруковано, і контур оцифровано з похибкою $\pm10\ldots\SI{20}{мм}$. Біля краю гофрування зростає приблизно в $\ee$ разів на кожні $R/N\approx\SI{0,13}{м}$, тому зсув осьової лінії на $\pm\SI{2}{см}$ змінює $\eripple(\SI{4,21}{м})$ до 4,18 і 3,22\,\% ($+14/-12\,\%$). При $R<\SI{4,1}{м}$ наші значення на 20--35\,\% нижчі за оцифровану криву джерела. Але там $\eripple<2\,\%$, і сама товщина лінії на рисунку ($\pm0{,}1\,\%$) порівнянна з різницею.
\end{projmodel}

\begin{figure}[!tb]\centering
  \includegraphics[width=\textwidth]{jet_ripple.jpg}
  \caption{Гофрування TF-поля JET 1975. Ліворуч --- на екваторі, логарифмічна шкала: наше $\dripple$ (суцільна), наше $\eripple=2\dripple$ (штрихова), точки --- значення, оцифровані нами з графіка гофрування~\cite[с.~330]{JET1975}, ромб --- 3,6\,\% на $R=\SI{4,21}{м}$~\cite[с.~83]{JET1975}. Праворуч --- $\eripple$ уздовж верхньої зовнішньої чверті межі: точки з таблиці~\cite[с.~332]{JET1975}, лінія --- наш розрахунок у тих самих точках. Власний розрахунок~\cite{ProjJET}.}
  \label{fig:04j-ripple}
\end{figure}

Чому гофрування так круто зростає назовні? Грубо кажучи, поле окремих ниток котушок, розташованих з кроком $2\pi R/N$ по колу радіуса $R_2$, спадає всередину як $(R/R_2)^N$, а від внутрішніх ніг на $R_1$ --- як $(R_1/R)^N$ (стандартна оцінка, поза корпусом). Для $N=32$, $R_1=\SI{1,30}{м}$, $R_2=\SI{4,79}{м}$ це дає $\dripple(\SI{4,21}{м})\approx(4{,}21/4{,}79)^{32}\approx1{,}6\,\%$ проти 1,83\,\% за Біо--Саваром. Для порівняння, у моделі ITER з розд.~\ref{ch:particles} на межі плазми $\dripple=4{,}5\cdot10^{-3}$. Плазма JET заповнює отвір майже до зовнішньої ноги, і за це доводиться платити гофруванням.

\section{Рівновага плазми в геометрії JET}\label{sec:04j-eq}

Щоб показати в реконструкції магнітні поверхні й силові лінії, потрібна рівновага в межах проєктної D-межі (розд.~\ref{ch:equilibrium}). Ми взяли найпростішу точну: розв'язок Соловйова.

\begin{outofcorpus}[рівновага Соловйова у формі Cerfon--Freidberg]
Якщо $\mu_0p'(\psi)$ і $FF'(\psi)$ сталі, рівняння $\GSop\psi=-\mu_0R^2p'-FF'$ стає лінійним. У змінних $x=R/R_0$, $y=Z/R_0$, $\psi=\psi_0u$ його розв'язок --- частинний поліном плюс сума семи однорідних розв'язків, парних по $y$. Сім коефіцієнтів визначають так, щоб межа $u=0$ проходила через три точки D-контуру (зовнішній і внутрішній екватор, верх) з правильними нахилом і кривинами~\cite{Cerfon2010}. Густина струму $\Jtor=(C_1R+C_2/R)/\mu_0$ від $\psi$ не залежить, тобто профіль \emph{плаский}.
\end{outofcorpus}

Межу задано у формі Міллера з $R_0$, $a$, $b$ проєкту та трикутністю $\delta=0{,}348$, підігнаною до форми межі з таблиці гофрування~\cite[с.~332]{JET1975}. Така плазма лежить у камері з зазорами 18 і \SI{5}{см}~\cite{ProjJET}. Сама таблиця суперечлива: її верхня точка ($Z=\SI{2,195}{м}$) на \SI{9,5}{см} вища за $b$ і виходить за стінку. Залишаються два вільні параметри. Поділ струму між $p'$ і $FF'$ вибрано так, щоб $\betap=0{,}9$, як у виносках проєкту, а масштаб $\psi$ --- так, щоб $\Ip=\SI{3,8}{МА}$~\cite[с.~83]{JET1975}. Результати --- у табл.~\ref{tab:04j-eq}.

\begin{table}[!tb]\centering\small
\caption{Рівновага Соловйова в межі JET 1975 (наш розрахунок~\cite{ProjJET}) і проєктні значення~\cite[с.~83]{JET1975}}\label{tab:04j-eq}
\begin{tabular}{@{}lll@{\hspace{2em}}lll@{}}
\toprule
Величина & Модель & Проєкт & Величина & Модель & Проєкт\\
\midrule
$\Ip$, МА & 3,80 & 3,8 & $\qnf$ & 5,51 & --- \\
$\Bzero$, Тл (вакуум) & 2,770 & 2,77 & $q(a)$ & 5,81 & 6 \\
$\betap$ & 0,90 & 0,9 & $q_0$ & 2,80 & $\approx1$ \\
$\betat$, \% & 2,32 & --- & $\li$ (3) / (1) & 0,425 / 0,498 & --- \\
$\Delta_{\mathrm S}$, м & 0,223 & --- & нев'язка межі, мм & 0,08 & --- \\
\bottomrule
\end{tabular}
\end{table}

Формула~\eqref{eq:04j-qD} дає $q(a)=5{,}99$, модель --- 5,81, тобто на 3\,\% менше. Автори проєкту самі зазначали, що навіть малі зміни форми перерізу змінюють $k$ і $q$ до 30\,\%~\cite[с.~83]{JET1975}. Зсув Шафранова можна порівняти з оцінкою з розд.~\ref{ch:equilibrium}: $(a^2/2R_0)(\betap+\li/2)\approx\SI{0,29}{м}$ для круглої плазми з великим аспектним відношенням. Витягнутість його зменшує, і модель дає \SI{0,223}{м}. \emph{Чесне обмеження.} Струм Соловйова плаский, тому $q_0=2{,}80$ і $q(a)/q_0=2{,}1$, а проєкт розраховував на пікований струм з $q(a)/q(0)=5{,}3\ldots6$, тобто $q_0\approx1$~\cite[с.~71, 83]{JET1975}. Тому поверхонь $q=1$, $3/2$ і $2$ у нашій плазмі \emph{немає}, і острів 2/1 тут нерезонансний. Модель придатна для топології (вкладені поверхні, раціональні лінії $q=3,\,7/2,\,4,\,5$ замикаються з точністю $<\SI{0,5}{мм}$), але не для МГД-стійкості. Для реалістичного $q$-профілю потрібен чисельний розв'язувач з $p'(\psi)$ і $FF'(\psi)$, що спадають до межі.

\section{Проєкт 1975 і машина 1983}\label{sec:04j-1983}

Між проєктом і першою плазмою минуло вісім років. Основні параметри плазми й магніту TF не змінилися: $a$, $b$, $R_0$, 32 котушки по 24 витки, \SI{12}{т} кожна, разом \SI{384}{т}~\cite[с.~13]{JETJU1983}, \cite[с.~6, 13]{JETBrochure1982}. Змінилося майже все довкола плазми (табл.~\ref{tab:04j-1983}, рис.~\ref{fig:04j-1983}).

\begin{table}[!tb]\centering\footnotesize\setlength{\tabcolsep}{4pt}\renewcommand{\arraystretch}{1.05}
\caption{Проєкт 1975~р. і машина 1983~р. (сторінки PDF)}\label{tab:04j-1983}
\begin{tabularx}{\textwidth}{@{}l>{\raggedright\arraybackslash}X>{\raggedright\arraybackslash}X@{}}
\toprule
Вузол & 1975~\cite{JET1975} & 1983~\cite{JETJU1983} та ін. \\
\midrule
Октант камери & 4 жорсткі сектори + 4 сильфони; разом 32 сектори (с.~295, 297) & 5 жорстких секторів трьох типів + 4 сильфони; разом 40 (с.~24); так само~\cite[с.~12]{JETBrochure1982} і вже в 1979~р.~\cite[с.~36]{JETJU1979} \\
Сильфони & лист 1,7~мм (с.~306) & лист 2~мм, гофри 120~мм (с.~24) \\
Маса / опір камери & 68~т; 0,52~мОм (с.~310) & 108~т (с.~13); 0,6~мОм (с.~24); кільця жорсткості на $Z=\pm\SI{1}{м}$ (с.~24) \\
Лімітери & 3 рейкові, 16 пластин Mo по $100\times60$~см (с.~301, 314) & 12 модулів по 0,32~м$^2$ на зовнішньому екваторі: 4 графітові + 8 з Ni-плакуванням на CuCr (с.~24); липень--листопад 1983~р. працювали лише 4 графітові (с.~39) \\
PF-1 & 10 котушок (с.~401) & 8 котушок, 40~кА; усього PF 14~\cite[с.~15]{JETBrochure1982} \\
Залізо & 2263~т (розширений варіант 2567~т) (с.~415) & 2800~т (с.~13) \\
Оболонка й кільця & зовнішня частина кільця --- литий легкий сплав (с.~376) & аустенітний високоміцний чавун, $\approx470$~т~\cite[с.~17]{JETBrochure1982} \\
Порти, нагрів & 4 насосні порти з 8 (с.~315); концепція NBI 6~інжекторів на порт (с.~426) & насосні камери на октантах 1 і 5, адаптери NBI на 4 і 8 (с.~24--25); поворотні клапани --- січень 1984~р., перший бокс NBI --- вересень 1984~р. (с.~40); у 1983~р. додаткового нагріву немає \\
Прогрів & до 500\,°C (с.~299) & проєктно 500\,°C; октанти відпалено при 520\,°C (с.~24); у кампаніях 1983~р. --- 270\,°C (с.~13), 300\,°C стінки і 200\,°C порти (с.~39) \\
\bottomrule
\end{tabularx}
\end{table}

\begin{figure}[!tb]\centering
  \begin{subfigure}{0.495\textwidth}\centering
    \includegraphics[width=\textwidth,trim=0 40 0 130,clip]{jet_invessel.jpg}\caption{}
  \end{subfigure}\hfill
  \begin{subfigure}{0.495\textwidth}\centering
    \includegraphics[width=\textwidth,trim=0 40 0 130,clip]{jet_cutaway1983.jpg}\caption{}
  \end{subfigure}
  \caption{(a)~Усередині камери JET 1975: стінки жорстких секторів (Inconel~600), гофровані сильфонні секції (Inconel~625), пластини зовнішнього рейкового Mo-лімітера, отвори портів. (b)~JET 1983 у тому самому ракурсі, що й рис.~\ref{fig:04j-cutaway}a: 40 жорстких частин камери, кільця жорсткості, дискретні модулі лімітера на зовнішньому екваторі, 8 котушок PF-1. Власна 3D-реконструкція~\cite{ProjJET}; побудовано за даними~\cite[с.~295--314]{JET1975} і~\cite[с.~13, 24--25, 39--40]{JETJU1983}.}
  \label{fig:04j-1983}
\end{figure}

Два пояснення до таблиці. \emph{40 жорстких частин при 32-кратній симетрії.} Уже в 1975~р. октанти зварювали посередині крайніх секторів, тож октант був послідовністю $\tfrac12+3+\tfrac12$ сектори~\cite[с.~299]{JET1975}. У 1983~р. ці половинки стали окремими деталями. Так ми тлумачимо «п'ять жорстких секторів» у~\cite[с.~24]{JETJU1983}; це тлумачення зберігає відповідність 32 котушкам TF. \emph{Графіт замість металу.} Проєкт лишав матеріал лімітера відкритим, і в машину поставили і графітові, і нікелеві модулі. Перші розряди 1983~р. були забруднені C і O; після підвищення температури прогріву плазма стала гарячішою, але з'явилися Ni і Cr зі стінки~\cite[с.~13]{JETJU1983}. Нікелеві модулі так і не випробували, а на 1984~р. планували пояскові лімітери й вивчали берилій~\cite[с.~39--40]{JETJU1983} (розд.~\ref{ch:wall}).

\paragraph{Суперечності джерел.} Їх варто знати, щоб не «виправляти» одне джерело іншим.
\begin{itemize}
\item \emph{Розміри камери.} «$4{,}3\times\SI{2,7}{м}$»~\cite[с.~23]{JETJU1983} --- внутрішня обшивка жорсткого сектора ($2\times2{,}135$ і $4{,}389-1{,}664$ за кресленням~\cite[с.~298]{JET1975}). «$2{,}5\times\SI{4,2}{м}$»~\cite[с.~12]{JETJU1983} --- це плазма $2a\times2b$. «$2{,}63\times\SI{4,18}{м}$»~\cite[с.~310]{JET1975}, \cite[с.~13]{JETBrochure1982} --- радіуси екранів сильфонів (4,33~м $-$ 1,70~м). Це три різні поверхні, а не три версії камери.
\item \emph{Маси.} Камера: 108~т~\cite[с.~13]{JETJU1983}, 100~т~\cite[с.~6]{JETBrochure1982}. Залізо: 2263 і 2567~т~\cite[с.~415]{JET1975}, 2567~т~\cite[с.~15]{JETJU1979}, 2700~т~\cite[с.~18]{JETBrochure1982}, 2800~т~\cite[с.~13]{JETJU1983}, 2600~т~\cite[с.~42]{Wesson1999}.
\item \emph{Температура лімітерів.} За звітом 1983~р. поверхня лімітера ніколи не нагрівалася понад 500\,°C~\cite[с.~39]{JETJU1983}, а за пізнішим оглядом типово сягала 800\,°C~\cite[с.~135]{Wesson1999}. Імовірно, мова про різні періоди роботи.
\end{itemize}

\begin{projmodel}[що в реконструкції оцифровано, а що припущено]\sloppy
\emph{З креслень, точно:} профіль камери. Ланцюг дуг з надрукованих радіусів замикається з похибкою $<\SI{0,2}{мм}$, товщина стінки 0,12--\SI{0,18}{м} збігається з таблицею~\cite[с.~298, 310]{JET1975}. Звідти ж PF-3 і PF-4 та кільце. \emph{Оцифровано:} D-форма котушки TF ($\pm10\ldots\SI{20}{мм}$, радіуси не надруковано), контур ярма ($\pm\SI{30}{мм}$), форма межі плазми. \emph{Припущено:} тороїдальна ширина плечей ярма ($\approx\SI{1,0}{м}$, виведено з площі перерізу \SI{11,2}{м^2} на 8 плечей~\cite[с.~412]{JET1975}), переріз PF-2 ($0{,}40\times\SI{0,40}{м}$, креслення немає), кількість і азимути вертикальних портів, положення кілець по $Z$ ($\pm\SI{0,15}{м}$), габарити інжекторів NBI ($\pm5\ldots\SI{10}{см}$). Для 1983~р. припущено кутову розкладку секторів, азимути 12 лімітерів і розміри насосних камер. Кожне число має в базі поле достовірності (\texttt{text/table/drawing/digitized/assumed}); див.\ \texttt{machines/jet1975/README.md} і \texttt{machines/jet1983/README.md}~\cite{ProjJET}.
\end{projmodel}

\begin{practicum}[JET у коді \texttt{tokamak-3d-viz}]\raggedright
З теки \texttt{tokamak-3d-viz} (інтерпретатор \texttt{PY="../fusion equilibrium challenge/starter/.venv/bin/python"}):
\begin{enumerate}
\item \texttt{\$PY src/run\_bake\_machine.py --machine machines/jet1975 --out data/bake\_jet1975}, потім те саме для \texttt{machines/jet1983} і \texttt{data/bake\_jet1983}. Далі \texttt{\$PY src/preview\_machine.py --bake data/bake\_jet1975}: створюються переріз $R$--$Z$ і план. Порівняйте \texttt{manifest.json} обох бейків: які компоненти зникли й які з'явилися, скільки записів позначено \texttt{assumed}.
\item \texttt{\$PY src/run\_bake\_jet\_physics.py} будує рівновагу, $q$, силові лінії, гофрування, \texttt{physics.json} і \texttt{ripple.png}. Очікувано: $\qnf=5{,}51$, $q(a)=5{,}81$, $\eripple(\SI{4,21}{м})=3{,}67\,\%$.
\item Тести: \texttt{\$PY -m pytest tests/test\_machine\_geometry.py tests/test\_machine\_1983.py tests/test\_jet\_physics.py} (42 тести).
\item \emph{Змініть число котушок з 32 на 16} за того самого струму (функції \texttt{tf\_filaments}, \texttt{ripple} у \texttt{src/tokviz/machine/ripple.py}; еталон --- \texttt{manual/code/ch04j\_ripple\_q.py}). Очікувано: $\dripple(\SI{4,21}{м})$ зростає з 1,83 до $\approx14{,}8\,\%$, а на осі --- з $\sim3\cdot10^{-7}$ до $\sim7\cdot10^{-4}$.
\item \emph{Змініть трикутність} (\texttt{build\_equilibrium(delta=...)} у \texttt{equilibrium\_analytic.py}). Очікувано $\qnf=4{,}81;\ 5{,}17;\ 5{,}51;\ 5{,}95$ при $\delta=0;\ 0{,}2;\ 0{,}348;\ 0{,}5$. Поясніть, чому $q$ зростає за сталого $\Ip$.
\end{enumerate}
Інтерактивний переглядач обох машин (вмикання вузлів, повзунок розрізу, підписи з розмірами) лежить у \texttt{tokamak-3d-viz/web}; самодостатню сторінку збирає \texttt{web/build\_page.py}.
\end{practicum}

\section*{Контрольні запитання}
\begin{enumerate}
\item Чому гнучка котушка в полі $\propto1/R$ набуває D-форми з $\rho\propto R$? Що відбувається з круговою котушкою тієї ж висоти і чому це небезпечно саме для епоксидної ізоляції?
\item Камера JET --- замкнений металевий виток навколо осердя трансформатора. Чому її тороїдальний опір роблять великим, а малий полоїдальний опір допустимий? Що змінилося б без сильфонів?
\item Проєкт наводить гофрування 3,6\,\%. Скільки це в сучасному означенні $\dripple$? Чому гофрування між $R_0$ і краєм плазми зростає на кілька порядків?
\item Навіщо залізне осердя, якщо центральна колона насичена? Що дає зворотне намагнічування перед імпульсом?
\item Чому рівновага Соловйова дає $q_0=2{,}8$, хоча проєкт закладав $q_0\approx1$? Яких раціональних поверхонь через це немає і які висновки з моделі не можна робити?
\end{enumerate}

\section*{Задачі}
\begin{enumerate}
\item \emph{Котушка TF.} Для JET 1975: $\ITF=\SI{41}{МА}$, 32 котушки по 24 витки, $R_0=\SI{2,96}{м}$, $L=\SI{0,66}{Гн}$, $I=\SI{53,4}{кА}$, осьова лінія від 1,304 до \SI{4,790}{м}, мідь внутрішньої ноги $24\times\SI{2700}{мм^2}$. (а)~Знайдіть $\Bzero$ і запасену енергію. (б)~За~\eqref{eq:04j-rho} знайдіть натяг $\Tcoil$ і напруження в міді; порівняйте з 5,9~кгс/мм$^2$~\cite[с.~339]{JET1975}. (в)~Перерахуйте $\Tcoil$ для \SI{51}{МА}.
\item \emph{Камера як виток.} $\Rv=\SI{0,52}{мОм}$, $\Lv=\SI{2,5}{мкГн}$ (1975) і $\SI{0,6}{мОм}$ (1983). (а)~Знайдіть $\Lv/\Rv$. (б)~У першій кампанії 1983~р. струм \SI{600}{кА} отримали при напрузі обходу \SI{14}{В}, а в листопаді \SI{1,9}{МА} --- при \SI{1,3}{В}~\cite[с.~13]{JETJU1983}. Який струм тече в камері і яку частку $\Ip$ він становить? Який опір плазми в обох випадках? (в)~Повторіть (б) для суцільної камери з опором \SI{0,06}{мОм} (оцінка з рамки). (г)~Оцініть розширення зовнішнього радіуса \SI{4,5}{м} при прогріві від 20 до \SI{500}{\degreeCelsius} (коефіцієнт лінійного розширення інконелю $1{,}4\cdot10^{-5}$~К$^{-1}$).
\enlargethispage{3\baselineskip}\item \emph{Форма і струм.} (а)~За оцінкою $(R/R_2)^N+(R_1/R)^N$ ($R_1=\SI{1,30}{м}$, $R_2=\SI{4,79}{м}$) знайдіть $\dripple$ і $\eripple$ на $R=\SI{4,21}{м}$ для $N=32$ і $N=16$ та порівняйте з Біо--Саваром (1,83 і 14,75\,\%). (б)~За першою з~\eqref{eq:04j-Ic} знайдіть $\ITF$ для круглої плазми ($\Ip=\SI{2,6}{МА}$, $q=2{,}8$, $R_0/a=2{,}37$). (в)~За~\eqref{eq:04j-qD} знайдіть $q(a)$ для $\Ip=\SI{3,8}{МА}$, $b/a=1{,}68$, $a/R_0=1{,}25/2{,}96$, $\ITF=\SI{41}{МА}$ і порівняйте з моделлю (5,81).
\end{enumerate}
